\documentclass[10pt,a4paper,twocolumn]{article}
% Shared preamble for the synthetic Turkish corpus (compiled with tectonic).
\usepackage{iftex}
\ifXeTeX
  \usepackage{fontspec} % Latin Modern OpenType, proper ToUnicode for Turkish letters
\else
  \usepackage[T1]{fontenc}
  \usepackage[utf8]{inputenc}
  \usepackage{lmodern}
\fi
\usepackage[turkish,shorthands=off]{babel}
\usepackage{amsmath,amssymb}
\usepackage{graphicx}
\usepackage{xcolor}
\usepackage{tikz}
\usepackage{booktabs}
\usepackage{hyperref}
\hypersetup{colorlinks=true,linkcolor=blue!50!black,urlcolor=blue!50!black,citecolor=green!40!black}
\newcommand{\paragrafA}{Makine çevirisi sistemleri son on yılda büyük bir gelişme göstermiştir. Özellikle sinir ağı tabanlı yöntemler, kelime düzeyinde değil tümce düzeyinde bağlamı dikkate alarak daha akıcı çeviriler üretmektedir. Bununla birlikte akademik metinlerde terimlerin tutarlı biçimde çevrilmesi, kesinlik derecesinin korunması ve sayısal değerlerin değiştirilmemesi hâlâ önemli sorunlar arasında yer almaktadır.}
\newcommand{\paragrafB}{Bu çalışmada 2019–2023 yılları arasında yayımlanmış 412 makaleden oluşan bir derlem kullanılmıştır. Her makale için özet, giriş ve sonuç bölümleri ayrı ayrı değerlendirilmiş; toplam 18.640 tümce iki bağımsız uzman tarafından incelenmiştir. Uzmanlar arasındaki uyum katsayısı $\kappa = 0{,}81$ olarak hesaplanmıştır.}
\newcommand{\paragrafC}{Bulgular, uzun ve iç içe geçmiş tümcelerde terim tutarlılığının belirgin biçimde azaldığını göstermektedir. Ortalama terim tutarlılığı \%87,4 olarak ölçülmüş; ancak kırk sözcükten uzun tümcelerde bu oran \%71,2'ye düşmüştür. Bu fark istatistiksel olarak anlamlıdır ($p < 0{,}01$).}

\usepackage[margin=1.8cm,columnsep=0.6cm]{geometry}
\begin{document}
\section{Bulgular}
Öğrencilerin tanıtıcı özellikleri iki üniversite için ayrı ayrı incelenmiştir.

A ve B üniversitelerinin son sınıf hemşirelik öğrencilerinin sosyodemografik özelliklerinin karşılaştırılması Tablo~1'de\newpage\noindent
sunulmuştur. A Üniversitesindeki öğrencilerin yaş ortalaması 22,4 olarak bulunmuştur.

\section{Kararlılık}
İspat. Fark şemasında $v_k^n = r^k e^{in\theta}$ yazılır ve Von-Neumann analiz metodu kullanılırsa
\begin{align}
\frac{r^2-2r+1}{\tau h^2} &= \frac{2r}{h^2}\left(\cos\theta - 1\right),\\
\frac{r^2-1}{2\tau} &= \frac{r}{h^2}\left(e^{i\theta} - 2 + e^{-i\theta}\right)
\end{align}
elde edilir. Burada $n=0$ ve $k=1$ alınırsa
\begin{align}
\frac{1}{\tau h^2}\left(r e^{-i\theta} - 2r^2\right) &= \frac{r}{h^2}\left(e^{i\theta} - 2\right),\\
\frac{r^2}{\tau} &= \frac{2r}{h^2}
\end{align}
olur. Bu denklemin kökleri birim çemberin içinde kaldığından şema kararlıdır.

\section{Öneriler}
\begin{itemize}
\item La charge épidémiologique des maladies bucco-dentaires est marquée par des inégalités affectant également les groupes spécifiques.
\item Binanın herhangi bir katında, her iki deprem doğrultusu için yapılmış olan performans analizi sonucu, kirişlerin en fazla yüzde yirmisi belirgin hasar bölgesine geçebilir.
\end{itemize}
\paragrafC
\end{document}
